\documentclass[10pt,a4paper]{amsart}
\usepackage[margin=19mm]{geometry}
\usepackage{amsmath,amssymb,mathtools}
\usepackage[scr=rsfs,cal=euler]{mathalfa}
\usepackage{xfrac}
\usepackage[unicode,hidelinks,bookmarks=false]{hyperref}
\newcommand{\ie}{i.e.\ }
\newcommand{\cf}{cf.\ }
\newcommand{\resp}{respectively\ }
\newcommand{\Subsection}{\subsection}
\providecommand{\nobreakdash}{\nobreak\hbox{-}}
\setlength{\emergencystretch}{2em}
\multlinegap0pt
\usepackage{enumitem}
\usepackage{needspace}
\usepackage{bookmark}
\setlength{\emergencystretch}{2em}
\numberwithin{equation}{section}
\newcommand{\F}{\mathbb F}
\newcommand{\PP}{\mathbb P}
\DeclareMathOperator{\Aut}{Aut}
\DeclareMathOperator{\GL}{GL}
\DeclareMathOperator{\PGL}{PGL}
\newtheorem{theorem}{Theorem}[section]
\newtheorem{proposition}[theorem]{Proposition}
\newtheorem{lemma}[theorem]{Lemma}
\newtheorem{corollary}[theorem]{Corollary}
\newtheorem*{theoremA}{Theorem~A}
\newtheorem{definition}[theorem]{Definition}
\newtheorem{remark}[theorem]{Remark}
\begin{document}
\title[A new family of maximal curves not covered by the Hermitian ]{A new family of maximal curves not covered by the Hermitian curve}
\author{Liming Ma; Yipeng Wang}
\date{}
\maketitle
\noindent\emph{Source and licence.} A new family of maximal curves not covered by the Hermitian curve. Version arXiv:2609.19546v1 (CC BY 4.0). This material is distributed under the Creative Commons Attribution 4.0 International licence, http://creativecommons.org/licenses/by/4.0/, as recorded in the arXiv OAI licence metadata for this exact version and shown on the abstract page https://arxiv.org/abs/2609.19546v1. Full rights notice: licenses/arxiv-2609.19546-CC-BY-4.0.txt.\par\smallskip
\noindent\textit{Editorial scope.} Counted unit: the original \texttt{theorem} environment \texttt{b1:noncover} at source lines 732--738 together with its complete proof at lines 740--828; the delivered document keeps source lines 116--870 so that every referenced label and every citation resolves inside it. Native editable source is the paper's own concatenated TeX; the AMS wrapper is editorial formatting only. No publisher class, upstream script or unrelated artwork is redistributed.
\begin{equation}
\label{b1:bm-model}
 x^{q+1}-y^{q+1}=1,\qquad w^m=\frac{x^{q^2}-x}{y^q}.
\end{equation}
We call these the GGS and BM families, respectively. Both are
maximal over \(\F_{q^{2n}}\) and have the same genus for each odd
\(n\). For \(n=3\), both are \(\F_{q^6}\)-isomorphic to the
GK curve; see \cite[Remark~2.1]{GGS2010} and
\cite[Section~2]{BM2018}.
For odd \(n\ge5\), they are not isomorphic over an algebraic closure
\cite[Corollary~2.6 and Remark~4.5]{BM2018}.

For odd \(n\ge5\), the GGS curves are not Galois subcovers of
\(\mathcal H_{q^n}\) over \(\F_{q^{2n}}\)
\cite{DM2012,GMZ2016}.
For \(q>2\) and odd \(n\ge3\), Beelen and Montanucci proved that the BM curves are
not Galois subcovers of
\(\mathcal H_{q^n}\) over \(\F_{q^{2n}}\)
\cite[Corollary~2.2]{BM2018}.
For odd \(n\ge5\), this result does not exclude non-Galois coverings.
The exponent \(m=(q^n+1)/(q+1)\) in the GGS and BM
constructions is integral only for odd \(n\).

In this paper we use a different Kummer model of the BM curves
which also defines maximal curves for even \(n\ge4\).
Let \(q\) be a prime power, let \(n\ge2\), and put
\(k=\F_{q^{2n}}\). Choose \(c\in k^*\) with \(c^{q^n}=-c\).
Let \(\mathcal X_{q,n}\) be the nonsingular projective curve over
\(k\) with function field \(k(u,z)\) given by
\begin{equation}
\label{b1:intro-model}
 z^{q^n+1}=c(u^q-u)
 \bigl(1+(u^q-u)^{q-1}\bigr)^{\frac{q^{n-1}-1}{q-1}}.
\end{equation}
Its \(k\)-isomorphism class is independent of the choice of \(c\).
For odd \(n\ge3\), an explicit change of generators gives a
\(k\)-isomorphism between \(\mathcal X_{q,n}\) and the
Beelen--Montanucci curve.
For \(n=2\), it is \(k\)-isomorphic to the Hermitian curve
\(\mathcal H_{q^2}\).
Our main results are summarized as follows.

\Needspace{10\baselineskip}
\begin{theoremA}
Let \(q\) be a prime power and let \(n\ge2\).
\begin{enumerate}
\item The curve \(\mathcal X_{q,n}\) is maximal over \(\F_{q^{2n}}\), and
its genus is
\[
 g(\mathcal X_{q,n})=
 \begin{cases}
 \bigl((q^2-1)q^n-q^2(q-1)\bigr)/2,&n\text{ odd},\\
 (q^2-1)q^n/2,&n\text{ even}.
 \end{cases}
\]
\item If \(q>2\) and \(n\ge3\), there is no nonconstant
\(\F_{q^{2n}}\)-morphism \(\mathcal H_{q^n}\to\mathcal X_{q,n}\).
\item If \(n\ge4\), then \(\Aut(\mathcal X_{q,n})\) contains a subgroup
of order \((q^n+1)q(q^2-1)\).
\end{enumerate}
\end{theoremA}

For \(q>2\) and odd \(n\ge5\), the second assertion strengthens
\cite[Corollary~2.2]{BM2018} by excluding arbitrary coverings.
For each fixed \(q>2\), taking \(n=2^{t-1}\) with \(t\ge3\)
gives an infinite family of \(\F_{q^{2^t}}\)-maximal curves
not covered by the Hermitian curve over the same field.

The restriction \(q>2\) is necessary: for \(q=2\), the curve
\(\mathcal X_{2,n}\) is a Galois subcover of \(\mathcal H_{2^n}\)
for every \(n\ge2\); see Remark~\ref{b1:q2-cover}.

We prove maximality by counting rational places that split completely
in \(k(u,z)/k(u)\), using fractional linear transformations over
\(\F_q\) and the norm from \(k\) to \(\F_{q^n}\).
The argument applies to both parities of \(n\) and gives an
independent proof of maximality for the BM curves.

To exclude a covering by the Hermitian curve, we represent the
pullbacks of a sequence of functions by homogeneous polynomials.
A degree bound turns their relations into polynomial identities,
and comparison of the multiplicities of an irreducible factor
gives a contradiction.

This paper is organized as follows. Section~\ref{b1:model-section} gives the change of generators from the
Beelen--Montanucci model and computes the genus.
Section~\ref{b1:maximality-section} proves maximality.
Section~\ref{b1:noncover-section} excludes coverings by the Hermitian
curve, and Section~\ref{b1:automorphisms-section} gives an explicit
subgroup of \(\Aut(\mathcal X_{q,n})\).


\section{A Kummer model and its genus}
\label{b1:model-section}

We retain the notation of the introduction and write \(p\) for the
characteristic of \(k\). We use the notation and terminology of
\cite{Stichtenoth2009} for algebraic function fields.
All isomorphisms and morphisms are over \(k\) unless another field
is specified.

Put \(A(T)=T^q-T\), \(B(T)=1+(T^q-T)^{q-1}\), and
\(r=(q^{n-1}-1)/(q-1)\).
Since \(B(T)=(T^{q^2}-T)/(T^q-T)\), its zeros are precisely the
elements of \(\F_{q^2}\setminus\F_q\), each with multiplicity one.
We write the function field of \(\mathcal X_{q,n}\) as
\begin{equation}
\label{b1:model}
 F=k(u,z),\qquad z^{q^n+1}=cA(u)B(u)^r.
\end{equation}

The right-hand side of \eqref{b1:model} has a simple zero at \(u=0\).
Eisenstein's criterion at \(u=0\) also applies over \(\overline k\).
Hence the defining polynomial is absolutely irreducible and
\([F:k(u)]=q^n+1\).
As \(q^n+1\mid q^{2n}-1\), all \((q^n+1)\)-st roots of unity lie in \(k\),
and \(F/k(u)\) is cyclic of degree \(q^n+1\).

The trace map from \(k\) to \(\F_{q^n}\) has a nonzero element in its
kernel, so the choice of \(c\) is possible.  If \(c'\) is another such
choice, then \(c'/c\in\F_{q^n}^*\).  Surjectivity of the norm
\(k^*\to\F_{q^n}^*\) gives \(\lambda\in k^*\) with
\(\lambda^{q^n+1}=c'/c\), and replacing \(z\) by \(\lambda z\)
gives the model with constant \(c'\).
Thus the \(k\)-isomorphism class of \(\mathcal X_{q,n}\) is independent
of the choice of \(c\).

\subsection{Relation with the Beelen--Montanucci model}

Let \(n\ge3\) be odd, and put \(m=(q^n+1)/(q+1)\).
Mendoza and Quoos gave a plane Kummer model for the BM curves
\cite[Corollary~3.3]{MQ2022}.
We give a change of generators from \eqref{b1:bm-model} to
\eqref{b1:model}, together with its inverse.

\begin{proposition}
\label{b1:bm-identification}
For odd \(n\ge3\), the curve \(\mathcal X_{q,n}\) is
\(\F_{q^{2n}}\)-isomorphic to the Beelen--Montanucci curve.
\end{proposition}

\begin{proof}
Let \(k(x,y,w)\) be the function field of the Beelen--Montanucci
curve given by \eqref{b1:bm-model}.
Choose \(\theta\in\F_{q^2}\setminus\F_q\). Since \(n\) is odd,
\(c=(\theta-\theta^q)^{-1}\) satisfies \(c^{q^n}=-c\).
The \(k\)-isomorphism class of \(\mathcal X_{q,n}\) is independent
of \(c\), so we use this choice in the proof.

The second equation in \eqref{b1:bm-model} can be written as
\(w^m=xy^{q^2}-yx^{q^2}\), since
\[
 y^q(xy^{q^2}-yx^{q^2})
 =x(x^{q+1}-1)^q-x^{q^2}(x^{q+1}-1)
 =x^{q^2}-x.
\]
Put \(s=x+y\) and \(t=\theta x+\theta^qy\). Then
\begin{equation}
\label{b1:bm-st}
 \begin{aligned}
 c(s^qt-st^q)&=x^{q+1}-y^{q+1}=1,\\
 c(s^{q^2}t-st^{q^2})&=xy^{q^2}-yx^{q^2}=w^m.
 \end{aligned}
\end{equation}
Setting \(u=s/t\), equation~\eqref{b1:bm-st} gives
\begin{equation}
\label{b1:bm-ab}
 A(u)=\frac{1}{ct^{q+1}},\qquad
 B(u)=\frac{u^{q^2}-u}{u^q-u}=\frac{w^m}{t^{q^2-q}}.
\end{equation}

Now put \(\ell=r/(q+1)=(q^{n-1}-1)/(q^2-1)\).
Since \(mr=\ell(q^n+1)\) and
\(q+1+(q^2-q)r=q^n+1\), equation~\eqref{b1:bm-ab} gives
\[
 cA(u)B(u)^r
 =\frac{w^{mr}}{t^{q+1+(q^2-q)r}}
 =\left(\frac{w^\ell}{t}\right)^{q^n+1}.
\]
Thus \eqref{b1:model} is satisfied by
\begin{equation}
\label{b1:bm-change}
 u=\frac{x+y}{\theta x+\theta^qy},\qquad
 z=\frac{w^\ell}{\theta x+\theta^qy}.
\end{equation}

\Needspace{10\baselineskip}
To recover the original generators, use
\(m-(q^2-q)\ell=1\) and \eqref{b1:bm-ab} to obtain
\[
 w=\frac{B(u)}{z^{q^2-q}},\qquad t=\frac{w^\ell}{z}.
\]
Then \(s=ut\), and solving for \(x,y\) gives
\[
 x=ct(1-\theta^qu),\qquad y=ct(\theta u-1).
\]
Thus \(k(x,y,w)=k(u,z)\), proving the assertion.
\end{proof}

\begin{remark}
\label{b1:small-members}
For \(n=2\), equation \eqref{b1:model} is
\(z^{q^2+1}=c(u^{q^2}-u)\).  With \(a=z\) and \(b=-cu\), this
becomes the Hermitian equation \(b^{q^2}+b=a^{q^2+1}\).
For \(n=3\), the curve \(\mathcal X_{q,3}\) is \(k\)-isomorphic to
the Giulietti--Korchm\'aros curve; see
Proposition~\ref{b1:bm-identification} and \cite[Section~2]{BM2018}.
\end{remark}

\subsection{Ramification and genus}

We compute the ramification in \(F/k(u)\). For \(a\in k\), let
\(Q_a\) denote the zero of \(u-a\) in \(k(u)\), and let
\(Q_\infty\) denote the pole of \(u\) in \(k(u)\).
The distinction between odd and even \(n\) occurs at the zeros of \(B\). Write
\(\varepsilon=\gcd(r,q^n+1)\).
The identity \(q^n+1=q(q-1)r+q+1\) gives
\(\gcd(r,q^n+1)=\gcd(r,q+1)\).  Reducing
\(r=1+q+\cdots+q^{n-2}\) modulo \(q+1\) gives zero for odd \(n\)
and one for even \(n\). Thus \(\varepsilon=q+1\) when \(n\) is odd,
and \(\varepsilon=1\) when \(n\) is even.

\begin{lemma}
\label{b1:ramification}
The following hold.
\begin{enumerate}
\item\label{b1:ramification-indices}
The only ramified places of \(k(u)\) in \(F\) are \(Q_\infty\)
and \(Q_a\) with \(a\in\F_{q^2}\). The \(q+1\) places
\(Q_\infty\) and \(Q_a\), \(a\in\F_q\), are totally ramified.
At the remaining \(q^2-q\)
places the ramification index is \((q^n+1)/(q+1)\) when \(n\) is
odd, and \(q^n+1\) when \(n\) is even.
\item\label{b1:infinite-place}
The unique place \(P_\infty\) of \(F\) lying above \(Q_\infty\)
is rational, and
\((u)_\infty=(q^n+1)P_\infty\), \((z)_\infty=q^nP_\infty\).
\end{enumerate}
\end{lemma}

\begin{proof}
In \(k(u)\), the function \(cA(u)B(u)^r\) has valuation one at
the places \(Q_a\) with \(a\in\F_q\), valuation \(r\) at the places
\(Q_a\) with \(a\in\F_{q^2}\setminus\F_q\), and valuation
\(-q^n\) at \(Q_\infty\). Since \(p\nmid q^n+1\),
\cite[Proposition~3.7.3]{Stichtenoth2009} gives the asserted
ramification indices and shows that every other place is unramified.
The place \(Q_\infty\) is rational and totally ramified, so \(P_\infty\) is
rational. The valuations of \(u\) and \(z\) at \(P_\infty\) follow
from \eqref{b1:model}.
\end{proof}

\begin{proposition}
\label{b1:genus}
The genus of \(\mathcal X_{q,n}\) is
\begin{equation}
\label{b1:genus-formula}
 g(\mathcal X_{q,n})=
 \begin{cases}
 \bigl((q^2-1)q^n-q^2(q-1)\bigr)/2,&n\text{ odd},\\
 (q^2-1)q^n/2,&n\text{ even}.
 \end{cases}
\end{equation}
\end{proposition}

\begin{proof}
Applying \cite[Corollary~3.7.4]{Stichtenoth2009} with the valuations
computed in the proof of Lemma~\ref{b1:ramification} gives
\[
 2g(F)-2=-2(q^n+1)+(q+1)q^n+(q^2-q)(q^n+1-\varepsilon).
\]
Substituting \(\varepsilon=q+1\) for odd \(n\) and
\(\varepsilon=1\) for even \(n\) gives
\eqref{b1:genus-formula}.
\end{proof}

For odd \(n\), \eqref{b1:genus-formula} agrees with the computation in
\cite[Proposition~2.1]{BM2018}.


\Needspace{18\baselineskip}
\section{Maximality}
\label{b1:maximality-section}

We prove maximality by counting rational places of \(k(u)\) that split
completely in \(F/k(u)\). Put \(f(T)=cA(T)B(T)^r\in k[T]\), so that
\eqref{b1:model} is \(z^{q^n+1}=f(u)\).
The map \(\lambda\mapsto\lambda^{q^n+1}\) is the norm map
from \(k^*\) onto \(\F_{q^n}^*\), with kernel of order \(q^n+1\).
Hence, for \(a\in k\setminus\F_{q^2}\), the place \(Q_a\) splits
completely in \(F/k(u)\) if and only if \(f(a)\in\F_{q^n}^*\).

\begin{lemma}
\label{b1:mobius-identities}
For \(M=\left(\begin{smallmatrix}\alpha&\beta\\\gamma&\delta\end{smallmatrix}\right)
\in\GL_2(q)\), write \(M(T)=(\alpha T+\beta)/(\gamma T+\delta)\).  Then
\begin{equation}
\label{b1:mobius-f}
 f(M(T))=\frac{\det(M)f(T)}{(\gamma T+\delta)^{q^n+1}}.
\end{equation}
Moreover, \(f(a)^{q^n}=-f(a^{q^n})\) for every \(a\in k\).
\end{lemma}

\begin{proof}
Using \(B(T)=(T^{q^2}-T)/A(T)\), substitution gives
\[
 A(M(T))=\frac{\det(M)A(T)}{(\gamma T+\delta)^{q+1}},\qquad
 B(M(T))=\frac{B(T)}{(\gamma T+\delta)^{q^2-q}}.
\]
Since \(q+1+(q^2-q)r=q^n+1\), these identities give
\eqref{b1:mobius-f}.  The last assertion follows from
\(A,B\in\F_q[T]\) and \(c^{q^n}=-c\).
\end{proof}

To find completely split places, we consider equations
\(a^{q^n}=M(a)\) with \(a\in k\setminus\F_{q^2}\) and
\(M\in\PGL_2(q)\). A trace-zero matrix representing \(M\) ensures
that such a solution satisfies \(f(a)^{q^n}=f(a)\), as shown in the
proof of Theorem~\ref{b1:maximality}.
Let \(\mathcal I_q\) be the set of elements of \(\PGL_2(q)\)
represented by matrices of trace zero in \(\GL_2(q)\).
For odd \(q\), these are precisely
the nonidentity involutions in \(\PGL_2(q)\). For even \(q\), scalar
matrices also have trace zero, so \(\mathcal I_q\) consists of the
nonidentity involutions and the identity.

\begin{lemma}
\label{b1:involution-solutions}
The following hold.
\begin{enumerate}
\item\label{b1:involution-cardinality}
The set \(\mathcal I_q\) has \(q^2\) elements.
\item\label{b1:involution-root-count}
For each
\(M\in\mathcal I_q\), the equation
\begin{equation}
\label{b1:frobenius-involution}
 a^{q^n}=M(a)
\end{equation}
has exactly \(q^n+1\) solutions in \(\PP^1(k)\).
\item\label{b1:involution-disjointness}
The solution sets of \eqref{b1:frobenius-involution} for
distinct elements of \(\mathcal I_q\) are disjoint outside
\(\PP^1(\F_{q^2})\).
\end{enumerate}
\end{lemma}

\begin{proof}
Each element of \(\mathcal I_q\) can be written uniquely in one of the forms
\[
 M(T)=-T+\beta\quad(\beta\in\F_q),\qquad
 M(T)=\frac{\alpha T+\beta}{T-\alpha}
 \quad(\alpha,\beta\in\F_q,\ \alpha^2+\beta\ne0).
\]
Hence \(|\mathcal I_q|=q+q(q-1)=q^2\).

In the first case, the finite solutions of
\eqref{b1:frobenius-involution} satisfy \(a^{q^n}+a=\beta\).
The trace map of \(k/\F_{q^n}\) gives \(q^n\) solutions in \(k\),
and \(\infty\) is also a solution.
In the second case, \eqref{b1:frobenius-involution} is equivalent to
\((a-\alpha)^{q^n+1}=\alpha^2+\beta\).
Since the right-hand side is nonzero, the norm map of
\(k/\F_{q^n}\) gives \(q^n+1\) solutions in \(k\).

If \(M(a)=M'(a)\) for distinct projective transformations over
\(\F_q\), then \(a\) satisfies a nonzero polynomial over \(\F_q\)
of degree at most two, or \(a=\infty\).  Thus
\(a\in\PP^1(\F_{q^2})\).
\end{proof}

\Needspace{8\baselineskip}
\begin{lemma}
\label{b1:ramified-rational-places}
The number of rational places of \(F\) lying above \(Q_a\),
\(a\in\PP^1(\F_{q^2})\), is \(q^3+1\) when \(n\) is
odd, and \(q^2+1\) when \(n\) is even.
\end{lemma}

\begin{proof}
For even \(n\), these places of \(k(u)\) are totally ramified in
\(F/k(u)\), by Lemma~\ref{b1:ramification}(\ref{b1:ramification-indices}), and give \(q^2+1\)
rational places of \(F\).

Suppose now that \(n\) is odd. The places \(Q_a\) with
\(a\in\PP^1(\F_q)\) give \(q+1\) rational places of \(F\).
Put \(m=(q^n+1)/(q+1)\). Since \(q+1\mid r\), the function
\(w=z^m/B(u)^{r/(q+1)}\) satisfies
\begin{equation}
\label{b1:intermediate-equation}
 w^{q+1}=cA(u).
\end{equation}
The function \(cA(u)\) has a simple zero at \(u=0\), so
\(T^{q+1}-cA(u)\) is irreducible and
\([k(u,w):k(u)]=q+1\).

For \(a\in\F_{q^2}\setminus\F_q\), we have
\(A(a)^{q^n}=-A(a)\) and \(c^{q^n}=-c\), hence
\(cA(a)\in\F_{q^n}^*\).  Since \((q+1)\mid q^n+1\), every element of
\(\F_{q^n}^*\) is a \((q+1)\)-st power in \(k\). Thus
\(T^{q+1}-cA(a)\) has \(q+1\) distinct roots in \(k\).
By \cite[Corollary~3.3.8(c)]{Stichtenoth2009}, the place \(Q_a\)
has \(q+1\) distinct rational extensions to \(k(u,w)\), and hence
splits completely.

Now \([F:k(u,w)]=m\). By
Lemma~\ref{b1:ramification}(\ref{b1:ramification-indices}), every place of \(F\) above
\(Q_a\) has ramification index \(m\) over \(k(u)\).
Since \(Q_a\) splits completely in \(k(u,w)\), multiplicativity of
ramification indices shows that each of these \(q+1\) rational
places of \(k(u,w)\) is totally ramified in \(F\).
Thus each has a unique rational place of \(F\) above it.
The total contribution for odd \(n\) is consequently \(q+1+(q^2-q)(q+1)=q^3+1\).
\end{proof}

\begin{theorem}
\label{b1:maximality}
For every prime power \(q\) and every integer \(n\ge2\), the curve
\(\mathcal X_{q,n}\) is maximal over \(\F_{q^{2n}}\).
\end{theorem}

\begin{proof}
We first consider the solutions of \eqref{b1:frobenius-involution}
in \(k\setminus\F_{q^2}\). Choose a trace-zero representative
\(M=\left(\begin{smallmatrix}\alpha&\beta\\\gamma&\delta\end{smallmatrix}\right)\)
and let \(a\in k\setminus\F_{q^2}\) satisfy this equation.
Using \(a^{q^n}=M(a)\) and \(\alpha+\delta=0\), we obtain
\[
 (\gamma a+\delta)^{q^n+1}
 =(\gamma M(a)+\delta)(\gamma a+\delta)
 =\gamma\beta+\delta^2=-\det(M).
\]
Lemma~\ref{b1:mobius-identities} now yields
\[
 f(a)^{q^n}=-f(a^{q^n})=-f(M(a))=f(a).
\]
Thus \(f(a)\in\F_{q^n}^*\), so the place \(Q_a\) splits completely
in \(F/k(u)\).

Suppose first that \(n\) is odd.  On \(\F_{q^2}\), the
\(q^n\)-power map is the \(q\)-power map. The trace and norm calculations
in Lemma~\ref{b1:involution-solutions}, applied to \(\F_{q^2}/\F_q\),
show that \(a^q=M(a)\) has exactly \(q+1\) solutions in
\(\PP^1(\F_{q^2})\). Therefore each element of \(\mathcal I_q\)
gives \(q^n-q\) rational places \(Q_a\), with \(a\notin\F_{q^2}\),
that split completely in \(F/k(u)\). The sets of places obtained
from distinct elements are disjoint by
Lemma~\ref{b1:involution-solutions}(\ref{b1:involution-disjointness}), giving
\(q^2(q^n-q)\) such places.

Suppose next that \(n\) is even.  The solutions in
\(\PP^1(\F_{q^2})\) are the fixed points of \(M\).
If \(q\) is odd, every involution has two distinct fixed points in
\(\PP^1(\F_{q^2})\).  Thus each of the \(q^2\) involutions contributes
\(q^n-1\) places.  If \(q\) is even, each nonidentity involution has
one fixed point, which belongs to \(\PP^1(\F_q)\), and contributes
\(q^n\) places.  The identity contributes \(q^n-q^2\) places.
The total in either characteristic is \(q^2(q^n-1)\).

Together with Lemma~\ref{b1:ramified-rational-places}, these
completely split places give
\begin{equation}
\label{b1:rational-count}
 \#\mathcal X_{q,n}(k)\ \ge\
 \begin{cases}
 (q^n+1)q^2(q^n-q)+q^3+1,&n\text{ odd},\\
 (q^n+1)q^2(q^n-1)+q^2+1,&n\text{ even}.
 \end{cases}
\end{equation}
By \eqref{b1:genus-formula}, each right-hand side equals
\(q^{2n}+1+2g(F)q^n\). The Hasse--Weil bound forces
equality, proving maximality.
\end{proof}

Equality in \eqref{b1:rational-count} shows that the places counted
in the proof are all the rational places of \(F\).
In particular, for \(a\in k\setminus\F_{q^2}\), the place \(Q_a\)
splits completely in \(F/k(u)\) if and only if
\eqref{b1:frobenius-involution} holds for some
\(M\in\mathcal I_q\); such an \(M\) is unique.

\begin{remark}
\label{b1:bm-maximality}
For odd \(n\ge3\), Theorem~\ref{b1:maximality} gives an independent
proof of the maximality established in \cite[Theorem~3.9]{BM2018}.
\end{remark}


\section{Coverings by the Hermitian curve}
\label{b1:noncover-section}

We prove that, for \(q>2\) and \(n\ge3\), the curve
\(\mathcal X_{q,n}\) is not covered by \(\mathcal H_{q^n}\) over
\(k\). Put \(N=q^n+1\) for the proof below. We first record a
sequence of functions on \(\mathcal X_{q,n}\) in the
Riemann--Roch space \(\mathcal L(NP_\infty)\).

For \(n\ge3\), put \(j=q^2-q\) and
\begin{equation}
\label{b1:ell}
 \ell=\left\lfloor\frac{r}{q+1}\right\rfloor
 =\begin{cases}
 (q^{n-1}-1)/(q^2-1),&n\text{ odd},\\
 (q^{n-1}-q)/(q^2-1),&n\text{ even},
 \end{cases}
 \qquad v=\frac{z^j}{B(u)}.
\end{equation}
Thus \(\ell\ge1\); when \(n\) is even, \(n\ge4\) and
\(\ell\ge q\).

\begin{lemma}
\label{b1:pole-functions}
For every prime power \(q\) and every \(n\ge3\), the function \(v\)
is nonconstant, and
\begin{equation}
\label{b1:function-list}
 1,u,z,zv,\ldots,zv^\ell\in\mathcal L(NP_\infty).
\end{equation}
\end{lemma}

\begin{proof}
By Lemma~\ref{b1:ramification}(\ref{b1:infinite-place}), the unique pole of \(u\) is
\(P_\infty\), and
\(\operatorname{ord}_{P_\infty}(u)=-N\),
\(\operatorname{ord}_{P_\infty}(z)=-q^n\).  Hence
\(\operatorname{ord}_{P_\infty}(v)=j\).
Put \(\varepsilon=\gcd(r,N)\). At every place of \(F\) above
\(Q_a\), with \(a\in\F_{q^2}\setminus\F_q\),
the orders of \(z\) and \(v\) are \(r/\varepsilon\) and
\(-(q+1)/\varepsilon\), respectively.  Since
\(r=(q+1)\ell\) when \(n\) is odd and
\(r=(q+1)\ell+1\) when \(n\) is even, the functions
\(zv^i\), \(0\le i\le\ell\), are regular at these places.
At every place of \(F\) above \(Q_a\), with \(a\in\F_q\), the orders of \(z\) and
\(v\) are \(1\) and \(j\), and there are no other finite poles.
This proves \eqref{b1:function-list}. The function \(v\) has a pole
above \(Q_a\) for every \(a\in\F_{q^2}\setminus\F_q\), so it is nonconstant.
\end{proof}

If a covering by \(\mathcal H_{q^n}\) exists, the following lemma
allows us to express the pullbacks of the functions in
\eqref{b1:function-list} as ratios of homogeneous polynomials.

\Needspace{8\baselineskip}
\begin{lemma}
\label{b1:polynomial-representation}
Let \(D\) be an effective \(k\)-rational divisor of degree
\(1\le d<q^n\) on \(\mathcal H_{q^n}:b^{q^n}+b=a^N\).
Let \(1,f_1,\ldots,f_s\) be nonzero functions in
\(\mathcal L(ND)\), computed over \(\overline{k}\).
\begin{enumerate}
\item\label{b1:homogeneous-representation}
There are homogeneous polynomials
\(F_0,\ldots,F_s\in\overline{k}[T_1,T_2,T_3]\) of degree \(d\)
such that
\(f_i=F_i(a,b,1)/F_0(a,b,1)\).
\item\label{b1:coprime-representation}
If \((f_1)_\infty=ND\), the polynomials in
\textup{(\ref{b1:homogeneous-representation})} can be chosen with
\(\gcd(F_0,F_1)=1\).
\item\label{b1:low-degree-nonvanishing}
No nonzero homogeneous polynomial of degree less than \(N\)
vanishes identically on \(\mathcal H_{q^n}\).
\end{enumerate}
\end{lemma}

\begin{proof}
Let \(O\) be the unique place at infinity of \(\mathcal H_{q^n}\).
By \cite[Corollary~1.2]{FGT1997},
\(q^nP+\operatorname{Fr}(P)\sim NO\) for every geometric point
\(P\), where \(\operatorname{Fr}\) is the \(q^{2n}\)-power
Frobenius. Summing over \(D\), and using
\(\operatorname{Fr}(D)=D\), gives
\begin{equation}
\label{b1:divisor-equivalence}
 ND\sim dNO.
\end{equation}
Choose a function \(h\) with \((h)=ND-dNO\). Then
\(h,hf_1,\ldots,hf_s\in\mathcal L(dNO)\).

The Weierstrass semigroup at \(O\) is \(\langle q^n,N\rangle\):
these are the pole orders of \(a,b\), and the semigroup they generate
has \(q^n(q^n-1)/2\) gaps, equal to the genus of \(\mathcal H_{q^n}\).
Every nongap at \(O\) has a unique expression \(iq^n+tN\) with
\(0\le i\le q^n\) and \(t\ge0\). Thus \(\mathcal L(dNO)\) has a basis consisting of the
monomials \(a^i b^t\) with \(0\le i\le q^n\), \(t\ge0\), and
\(iq^n+tN\le dN\): their pole orders are pairwise distinct and
exhaust the nongaps not exceeding \(dN\).
Since \(d<q^n\), the inequality \(i+t>d\) would imply
\(iq^n+tN\ge(d+1)q^n>dN\). Homogenizing the resulting expressions
for \(h,hf_1,\ldots,hf_s\) therefore gives the required polynomials.
This proves \textup{(\ref{b1:homogeneous-representation})}.

For \textup{(\ref{b1:low-degree-nonvanishing})}, let \(G\) be a nonzero homogeneous
polynomial with \(\deg G<N\). The monomials in \(G(a,b,1)\) have
distinct pole orders at \(O\), so \(G(a,b,1)\ne0\).

To prove \textup{(\ref{b1:coprime-representation})}, put
\(E_G=(G(a,b,1))+N(\deg G)O\) for each such polynomial \(G\).
Each monomial has pole order at most \(N\deg G\), so \(E_G\) is
effective of degree \(N\deg G\).
These divisors satisfy \(E_{GH}=E_G+E_H\)
whenever \(\deg(GH)<N\).
If \((f_1)_\infty=ND\), the choice of \(h\) and the homogenization give
\[
 \begin{aligned}
 E_{F_0}&=(h)+dNO=ND=(f_1)_\infty,\\
 E_{F_1}&=(hf_1)+dNO=(f_1)_0.
 \end{aligned}
\]
These two divisors have disjoint supports. A common factor
\(G\) of positive degree would give a nonzero effective divisor
\(E_G\) contained in both, a contradiction.
\end{proof}

\Needspace{6\baselineskip}

\begin{theorem}
\label{b1:noncover}
Let \(q>2\) and \(n\ge3\). The curve \(\mathcal X_{q,n}\) is not
covered by the Hermitian curve \(\mathcal H_{q^n}\) over \(\F_{q^{2n}}\);
that is, there is no nonconstant \(\F_{q^{2n}}\)-morphism
\(\mathcal H_{q^n}\to\mathcal X_{q,n}\).
\end{theorem}

\begin{proof}
Suppose first that there is a separable \(k\)-morphism
\(\pi:\mathcal H_{q^n}\to\mathcal X_{q,n}\) of degree \(d\).
The Hurwitz formula gives
\begin{equation}
\label{b1:hurwitz-bound}
 d(2g(\mathcal X_{q,n})-2)\le q^{2n}-q^n-2.
\end{equation}
Equation \eqref{b1:genus-formula} gives
\[
 2g(\mathcal X_{q,n})-2-(j+1)(q^n-2)
 \ge(q-2)(q^n-q^2+q)>0.
\]
Together with \eqref{b1:hurwitz-bound}, this yields
\begin{equation}
\label{b1:polynomial-degree-bound}
 (j+1)d<N.
\end{equation}
In particular, \(d<q^n\).
Since \(N=jr+q+1\), \(r\ge q+1\), and
\(r\le(q+1)\ell+1\), we obtain
\begin{equation}
\label{b1:cover-degree}
 d<\frac{N}{j+1}\le r\le(q+1)\ell+1<j\ell+1,
\end{equation}
where the last inequality uses \(q>2\), hence \(j>q+1\), and
\(\ell\ge1\).

Put \(D=\pi^*P_\infty\). This is an effective \(k\)-rational
divisor of degree \(d\). The pullbacks of \eqref{b1:function-list}
belong to \(\mathcal L(ND)\), and the pullback of \(u\) has pole
divisor \(ND\). By Lemma~\ref{b1:polynomial-representation}(\ref{b1:homogeneous-representation})--(\ref{b1:coprime-representation}), they
are represented by nonzero homogeneous polynomials
\(X_0,X_1,Y_0,\ldots,Y_\ell\) of degree \(d\) over \(\overline{k}\),
with \(\gcd(X_0,X_1)=1\). Their restrictions to \(\mathcal H_{q^n}\)
satisfy \(u=X_1/X_0\) and \(zv^i=Y_i/X_0\), where we suppress
the pullback notation.

On \(\mathcal H_{q^n}\), the polynomials satisfy
\begin{equation}
\label{b1:quadratic-identities}
 Y_0Y_{i+1}=Y_1Y_i,\qquad 0\le i<\ell.
\end{equation}
Since
\(B(u)=\prod_{\alpha\in\F_{q^2}\setminus\F_q}(u-\alpha)\),
the relation \(vB(u)=z^j\) also gives
\begin{equation}
\label{b1:factor-identity}
 Y_1\prod_{\alpha\in\F_{q^2}\setminus\F_q}
 (X_1-\alpha X_0)=Y_0^{j+1}.
\end{equation}
The homogeneous equations \eqref{b1:quadratic-identities} and
\eqref{b1:factor-identity} have degrees \(2d\) and \((j+1)d\),
respectively. Both are less than \(N\) by
\eqref{b1:polynomial-degree-bound}. Hence
Lemma~\ref{b1:polynomial-representation}(\ref{b1:low-degree-nonvanishing}) shows that these equations
are polynomial identities.

The restriction of \(Y_1/Y_0\) is the pullback of the nonconstant
function \(v\). As \(Y_0,Y_1\) have the same degree,
\(Y_0\) does not divide \(Y_1\). Choose an irreducible factor
\(\rho\) whose multiplicities in \(Y_0,Y_1\) are \(a_0\) and
\(a_0-e\), respectively, with \(e>0\).
By \eqref{b1:quadratic-identities}, the multiplicity of \(\rho\)
in \(Y_i\) is \(a_0-ie\). Since the multiplicity \(a_0-\ell e\) of
\(\rho\) in \(Y_\ell\) is nonnegative, \(a_0\ge\ell e\).
By \eqref{b1:factor-identity}, the multiplicity of \(\rho\) in
\(\prod_\alpha(X_1-\alpha X_0)\) is
\(ja_0+e\ge(j\ell+1)e>d\), where the last inequality follows
from \eqref{b1:cover-degree}.
However, \(\gcd(X_0,X_1)=1\) makes the factors
\(X_1-\alpha X_0\) pairwise relatively prime.
Thus \(\rho\) divides at most one of them and has multiplicity
at most \(d\), a contradiction.

Finally, suppose that an inseparable \(k\)-morphism exists, and
view \(F\) as a subfield of \(k(a,b)\) through its pullback.
Let \(E\) be the maximal intermediate field separable over \(F\).
Then \(k(a,b)/E\) is purely inseparable of degree \(p^s\) for
some \(s\ge1\). Since \(k\) is perfect,
\cite[Proposition~3.10.2(c)]{Stichtenoth2009} gives
\(E=k(a,b)^{p^s}=k(a^{p^s},b^{p^s})\).
The equation \(b^{q^n}+b=a^{q^n+1}\) has
coefficients in \(\F_p\), so the substitutions
\(a\mapsto a^{p^s}\), \(b\mapsto b^{p^s}\), fixing \(k\),
give a \(k\)-isomorphism from \(k(a,b)\) onto \(E\).
This would give a separable \(k\)-morphism
\(\mathcal H_{q^n}\to\mathcal X_{q,n}\), a contradiction.
\end{proof}



\begin{remark}
\label{b1:q2-cover}
For \(q=2\), the curve \(\mathcal X_{2,n}\) is a cyclic Galois
subcover of \(\mathcal H_{2^n}\) for every \(n\ge2\).
For odd \(n\), a cyclic Galois covering of degree \((2^n+1)/3\)
follows from Proposition~\ref{b1:bm-identification} and the proof of
\cite[Lemma~2.4]{BM2018}.

Suppose that \(n\) is even. We take \(c=1\), since the
\(k\)-isomorphism class is independent of \(c\), and choose
\(\omega\in\F_4\setminus\F_2\). Put \(d=(2^n-1)/3\) and use
the model \(y^{2^n+1}=x^{2^n}+x\) for \(\mathcal H_{2^n}\).
A covering of degree \(d\) is given by
\begin{equation}
\label{b1:q2-even-map}
 u=\frac{x^d+1}{\omega x^d+\omega^2},\qquad
 z=\frac{yx^{(d-1)/2}}{\omega x^d+\omega^2}.
\end{equation}
Put \(t=x^d\). Using \(y^{2^n+1}=x^{2^n}+x\),
\(3d=2^n-1\), and \(r=2^{n-1}-1\), we obtain
\[
 z^{2^n+1}=\frac{t^r(t^3+1)}{(\omega t+\omega^2)^{2^n+1}}
 =A(u)B(u)^r.
\]

Since \(k(u)=k(t)\) and \(k(x,y)=k(u,z)(x)\), the relation
\(x^d=t\) shows that the covering has degree at most \(d\).
As \(d\mid2^n-1\), the maps
\((x,y)\mapsto(\zeta^2x,\zeta y)\), with
\(\zeta\in k\) and \(\zeta^d=1\), form a cyclic group of order
\(d\) fixing \(u,z\). Hence the covering is cyclic Galois of
degree \(d\).
For \(n=2\), formula~\eqref{b1:q2-even-map} gives an isomorphism.
\end{remark}


\Needspace{10\baselineskip}
\section{An explicit automorphism subgroup}
\label{b1:automorphisms-section}


\begin{thebibliography}{99}
\bibitem[BM2018]{BM2018} P.~Beelen and M.~Montanucci, \href{https://doi.org/10.1112/jlms.12144}{\emph{A new family of maximal curves}}, J. London Math. Soc. (2) \textbf{98} (2018), no.~3, 573--592.
\bibitem[DM2012]{DM2012} I.~Duursma and K.-H.~Mak, \href{https://doi.org/10.1007/s00574-012-0022-2}{\emph{On maximal curves which are not Galois subcovers of the Hermitian curve}}, Bull. Braz. Math. Soc. (N.S.) \textbf{43} (2012), no.~3, 453--465.
\bibitem[FGT1997]{FGT1997} R.~Fuhrmann, A.~Garcia, and F.~Torres, \href{https://arxiv.org/abs/alg-geom/9610023}{\emph{On maximal curves}}, J. Number Theory \textbf{67} (1997), no.~1, 29--51.
\bibitem[GGS2010]{GGS2010} A.~Garcia, C.~G\"uneri, and H.~Stichtenoth, \href{https://doi.org/10.1515/ADVGEOM.2010.020}{\emph{A generalization of the Giulietti--Korchm\'aros maximal curve}}, Adv. Geom. \textbf{10} (2010), no.~3, 427--434.
\bibitem[GMZ2016]{GMZ2016} M.~Giulietti, M.~Montanucci, and G.~Zini, \href{https://doi.org/10.1016/j.ffa.2016.05.005}{\emph{On maximal curves that are not quotients of the Hermitian curve}}, Finite Fields Appl. \textbf{41} (2016), 72--88.
\bibitem[MQ2022]{MQ2022} E.~A.~R.~Mendoza and L.~Quoos, \href{https://doi.org/10.1016/j.ffa.2021.101945}{\emph{Explicit equations for maximal curves as subcovers of the BM curve}}, Finite Fields Appl. \textbf{77} (2022), Art.~101945, 22~pp.
\bibitem[Stichtenoth2009]{Stichtenoth2009} H.~Stichtenoth, \href{https://doi.org/10.1007/978-3-540-76878-4}{\emph{Algebraic Function Fields and Codes}}, second edition, Graduate Texts in Mathematics \textbf{254}, Springer, Berlin, 2009.
\end{thebibliography}
\end{document}
